\documentclass[a4paper,12pt]{article}
\usepackage[brazilian]{babel}
\usepackage[utf8]{inputenc}
\usepackage[T1]{fontenc}
\usepackage{indentfirst}
\usepackage{graphicx}
\usepackage{fancyhdr}
\usepackage{setspace}
\usepackage[a4paper,left=3cm,right=3cm,top=3cm,bottom=3cm]{geometry}
\usepackage{times}
% include --------------------------------------------------------------------------------------------------------------------
% \pagestyle{fancy}
\lhead{
  \begin{minipage}[c]{0.1\linewidth}
   \includegraphics[scale=0.5]{images/logoUnB_mini.png}
  \end{minipage}
}
\rhead{
  \hfill
  \begin{minipage}[c]{0.9\linewidth}
   \fontsize{6}{8}\selectfont
   Universidade de Brasília \\
   Instituto de Exatas - IE \\
   Programa de Pós-gradução em Informática \\
   Edital n. 02/2015
  \end{minipage}
}
\cfoot{
\fontsize{6}{8}\selectfont
\textbf{Universidade de Brasília – Campus Universitário Darcy Ribeiro – CEP 70910-900 – Brasília/DF - Brasil}\\
Instituto de Ciências Exatas - Departamento de Ciência da Computação – Caixa postal: 4466\\
Telefone/Fax: (61) 3107-3661\\
www.cic.unb.br / ppgi.unb.br/
}

% include --------------------------------------------------------------------------------------------------------------------

\begin{document}
\onehalfspacing
%comandos-----------------------------------------------------------------------------------------------------------------
\newcommand{\linhaDupla}{\newline \newline}

%comandos-----------------------------------------------------------------------------------------------------------------

% \begin{figure}
%  \includegraphics[width=0.8\textwidth]{images/logoUnB.png}
% \end{figure}
% 
% \textbf{PROGRAMA DE PÓS-GRADUAÇÃO EM INFORMÁTICA}
% \linhaDupla
% \textbf{Nome:} Leonardo Henrique Moreira
% \linhaDupla
% \textbf{CPF:} 294.490.948-70
% \linhaDupla
% \textbf{Proposta de Projeto de Pesquisa}
% \linhaDupla
% \textbf{Título:}
% \linhaDupla
% \textbf{Linha de Pesquisa:}
% \linhaDupla
% \indent\textbf{Área de Pesquisa:}
% 
% \newpage
\title{Proposta de emprego de ordenação de informações como ferramenta do processo de apoio à decisão}
\maketitle
\section{Introdução}
A capacidade do homem em perceber, interpretar e interagir com o ambiente no qual está inserido depende de um elemento essencial, a informação.
O estado atual da tecnologia, aliado à sua capacidade de cobertura global e em tempo real, aumenta a facilidade de acesso à informação e dessa forma as pessoas 
tem sido apresentadas a um volume de dados cada vez maior. Apesar de muitos destes dados serem relevantes, existem também aqueles que não agregam valor.\\
\indent De maneira especial, em um cenário militar, essa grande quantidade de informação pode prejudicar o processo decisório do 
comandante pois este volume, aliado à velocidade com a qual elas são apresentadas, pode dificultar a definição de qual tratamento cada uma deve receber.
A sobrecarga do decisor pode causar tanto atenção excessiva a fatos pouco relevantes quanto distração sobre pontos importantes.\\
\indent A condição ideal para o decisor, em qualquer escalão, é que exista um equilíbrio entre o fluxo de informação e o processo de interpretação dos dados.
Entretanto, apesar da alta taxa na qual dados tem sido coletados oriundos do grande número de sensores e sistemas utilizados para os mais diversos fins, 
verifica-se que a presença de ferramentas e algoritmos capazes de tratar esse volume de informação não existem \cite{Savas 2014}.\\
\indent Face a este problema, esta proposta de estudo tem como objetivo construir um modelo capaz de ordenar as informações 
de tal forma que o valor final agregado ao decisor seja maximizado. Como consequência dessa contrução, será possível verificar 
tanto se é possível associá-lo a sistemas de apoio à decisão, à teoria dos jogos e à técnicas de inteligência artificial quanto 
avaliar o ganho que o modelo proporciona.
\indent A construção deste modelo é importante porque através dele é possível auxiliar o processo de tomada de decisão devido ao fato de evitar ou diminuir 
a sobrecarga de informação dos decisores. Outro aspecto que aumenta a contribuição desta proposta de trabalho é o contínuo emprego das Forças Armadas e 
Auxiliares (Corpo de Bombeiros e Polícia Militar). No contexto do Exército Brasileiro, este trabalho auxiliará o exercício do Comando e Controle em operações 
de garantia da lei e da ordem e de defesa interna, atividades nas quais a Força Terrestre tem sido empregada frequentemente. No cenário das Forças Auxiliares,
a contribuição será na agilidade na priorização das demandas diárias.

\section{Justificativa}
A doutrina de Comando e Controle \cite{Brasil} relaciona a qualidade de informação por meio de sete critérios: precisão, oportunidade, facilidade de uso,
suficiência, brevidade e segurança. O segundo item merece atenção especial por relacionar o instante da coleta da informação com o momento até o qual ela é 
útil para o processo de toma de decisão.\\
\indent Segundo \cite{Savas 2014}, não existem ferramentas e algoritmos capazes de tratar o volume de informação, principalmente nos critérios supracitados. Aliando à
essa condição a busca pela superioridade de informação, isto é, a habilidade de coletar e processar dados, torna-se evidente a necessidade de
contribuições científicas e técnicas nesta área.

\section{Objetivos}
\indent O objetivo do estudo é criar um método capaz de realizar o ordenamento da lista de informações presente no sistema de 
Comando e Controle utilizado. A motivação para tal tratamento é oferecer ao decisor a possibilidade de avaliar, de maneira 
prioritária, as informações mais importantes ou mais urgentes, uma vez que tratá-las na ordem de chegada pode não ser 
eficiente.\\
\indent Como o exercício de C2 está presente nos diversos níveis hierárquicos das operações, a valorização de cada item 
da lista depende da sua função de combate (movimento e manobra, inteligência, logística, fogos, proteção), do escalão de 
seu decisor, da urgência no seu atendimento e principalmente do seu instante de chegada. Outra característica importante 
deste método deve ter é ser capaz de ser implementado e integrado à Família de Aplicativos de Comando e Controle da 
Força Terrestre.\\
\section{Revisão da literatura}
Os desafios no cenário de Comando e Controle são apresentados em \cite{Athans 1986}, onde são destacadas as necessidades de avanços na área da teoria 
da decisão e de estudo do impacto do atraso no tratamento de informação.\\
\indent O Modelo Dinâmico de Cognição Situada (\textit{Dynamic Model of Situated Cognition - DMSC}) foi introduzido pela primeira vez em 2003 e rapidamente 
ganhou aceitação na comunidade do DoD (\textit {Departament of Defense}) como um modelo descritivo de processamento e fluxo de dados e informacoes no campo de 
batalha \cite{Phillips 2008}.\\
\indent Este modelo demonstra a relação entre os sistemas tecnológicos e processos perceptivos e cognitivos humanos, fornecendo um quadro eficaz no qual se pode
ver o fluxo de dados e informações em um sistema complexo. Apesar de não ser um modelo analítico, o DMSC fornece um meio de ver o problema do fluxo de dados e 
informações de uma forma que seja acessível não somente para os analistas como também para os decisores. O DMSC é a base para a construção do modelo matemático
 para qualificação da informação.\\
\indent Além desse trabalho, outras pesquisas relatam a análise de informação no cenário de Comando e Controle. Entre elas destacam-se o uso de teoria dos 
jogos para predição de ameaças e análise de situação \cite{Brynielsson 2004}, na qual é estudada a aplicação de teoria de jogos como ferramenta de apoio à 
decisão favorecendo a construção da consciência situacional.


\section{Metodologia}
\indent Para atingir o objetivo estabelicido, será primeiramente conduzida uma pesquisa bibliográfica relativa a Comando e Controle, buscando mapear as ligações 
com as áreas de inteligência artificial, teoria dos jogos e sistema de apoio à decisão. A partir daí, será feito um estudo mais 
detalhado do processo de qualificação da informação com objetivo de produzir um modelo matemático que considere os critérios, principalmente o da oportunidade.\\
\indent Com a definição desse modelo, serão estudados cenários de uso e serão realizadas simulações verificando qual o ganho do ordenamento das informações para
 o processo decisório.
\section{Plano de trabalho}
O curso será realizado através das seguintes atividades (não necessariamente sequenciais):
\begin{itemize}
 \item cursar disciplinas necessárias;
 \item executar pesquisa bibliográfica sobre Comando e Controle, inteligência artificial, teoria dos jogos e sistemas de apoio 
 à decisão;
 \item defender qualificação com proposta de pesquisa específica;
 \item construir modelo matemático de qualificação da informação;
 \item estudar de cenários de uso e realizar simulações com o modelo proposto;
 \item analisar desempenho do modelo matemático;
  \item elaborar a dissertação;
 \item defender a dissertação;
\end{itemize}

% \newpage
\section{Cronograma}
\begin{figure}[!hbp]
  \centering
  \includegraphics[scale=0.5]{cronograma.png}
\end{figure}
% \section{Referência bibliográficas}

\begin{thebibliography}{1}

\bibitem{Brasil}
Brasil, \emph{MD31-D-03: Doutrina Militar de Comando e Controle}, Brasília, 2006.
\bibitem{Phillips 2008}
D. D. Phillips, \emph{Mathematical Modeling and Optimal Control of Battlefield Information Flow}, Naval Postgraduate School, 
2008.

\bibitem{Alberts 2006}
D. S. Albert, R. E. Hayes, \emph{Understanding Command and Control}, CCRP, 2006.

\bibitem{Athans 1986}
M. Athans, \emph{Command and Control (C2) Theory: A challenge to control science}, 1986.

\bibitem{Brynielsson 2004}
J. Brynielsson, S. Arnborg, \emph{Bayesian Games for Threat Prediction and Situation Analysis}, Seventh International
Conference on Information Fusion, Sweden, 2004.

\bibitem{Savas 2014}
O. Savasa, Y. Sadguyu, J. Deng, J. Li, \emph{Tactical big data analytics: challenges, use cases and solutions.} In 
\textit{ACM SIGMETRICS Performance Evaluation Review}, volume 41, pages 86-89, 2014.

\end{thebibliography}

\end{document}
